\chapter{Anexos}
\label{sec.anexos:codigo}

\section{Código Fuente del Ejercicio Práctico N.º 1}

A continuación se presenta el código completo del archivo \texttt{main.c} correspondiente al Ejercicio 1,
donde se implementan las tareas de señalización luminosa y telemetría junto al gancho de inactividad.

\lstinputlisting[
  language=C,
  caption={Código fuente completo de \texttt{main.c} para el Ejercicio 1.},
  label={list.ej1:main}
]{sources/ep1/src/main.c}

\newpage

\section{Código Fuente del Ejercicio Práctico N.º 2}

A continuación se transcribe el código del archivo \texttt{main.c} correspondiente al Ejercicio 2,
incorporando el periférico ADC1, la rutina de interrupción \texttt{ADC1\_2\_IRQHandler}, la cola \texttt{xADC1Queue}
y el semáforo binario \texttt{xADC1Semaphore}.

\lstinputlisting[
  language=C,
  caption={Código fuente completo de \texttt{main.c} para el Ejercicio 2.},
  label={list.ej2:main}
]{sources/ep2/src/main.c}

\newpage

\section{Configuración del Kernel (\texttt{FreeRTOSConfig.h})}

En este anexo se adjunta el archivo de cabecera \texttt{FreeRTOSConfig.h} empleado para parametrizar el núcleo del
sistema operativo en ambos ejercicios prácticos.

\lstinputlisting[
  language=C,
  caption={Archivo de configuración \texttt{FreeRTOSConfig.h} del kernel para STM32F103.},
  label={list.config:freertos}
]{sources/ep1/include/FreeRTOSConfig.h}
